% ============================================================================
% Complexity-Aware Dynamic Precision Routing for Energy-Proportional
% LLM Inference — IEEEtran conference draft
%
% EDITORIAL NOTE: This is a LaTeX rendering of major-project/paper/draft.md,
% produced as a direct 1:1 conversion for IEEE submission typesetting. It
% carries the same evidentiary constraints as the markdown source: every
% number in this file must trace to a PASS row in major-project/paper/
% results.md; RQ2/RQ3/RQ4 are stated as pending with zero projected outcome.
% Author names/affiliations and the bibliography below are placeholders —
% see the \TODO markers — and must be completed/verified before submission.
%
% NOTE ON PROVENANCE (2026-09-08): the Results section (Sec. V-A) reflects
% the Session 1 numbers as recorded in paper/results.md on the local `main`
% branch at the time this .tex was generated. A parallel, more advanced
% line of work exists on origin/main (10 commits ahead, including partial
% Session 4 routing-experiment data and an active investigation into why
% the fuzzy router underperforms a matched-random baseline) that has not
% yet been reconciled into this draft. Do not treat Sec. V-A/B/C/D as final
% until that reconciliation happens — see the project roadmap (NEW.md) and
% CREDIBILITY_REPORT.md for current status before this is submitted anywhere.
% ============================================================================

\documentclass[conference]{IEEEtran}
\IEEEoverridecommandlockouts

\usepackage{cite}
\usepackage{amsmath,amssymb,amsfonts}
\usepackage{algorithmic}
\usepackage{graphicx}
\usepackage{textcomp}
\usepackage{xcolor}
\usepackage{booktabs}
\usepackage{hyperref}

\def\BibTeX{{\rm B\kern-.05em{\sc i\kern-.025em b}\kern-.08em
    T\kern-.1667em\lower.7ex\hbox{E}\kern-.125emX}}

\begin{document}

\title{Complexity-Aware Dynamic Precision Routing for\\Energy-Proportional LLM Inference\\
\large\textit{Working title --- markdown-to-LaTeX working draft}}

\author{
\IEEEauthorblockN{Author names withheld}
\IEEEauthorblockA{\textit{Affiliation to be added at submission} \\
\textit{Green-Weight Project}\\
email@TODO}
}

\maketitle

\begin{abstract}
Large language model (LLM) inference is typically served at a single,
statically chosen numeric precision for an entire deployment. This paper
studies an alternative: routing each individual request, at inference
time, to one of several precision tiers of the \emph{same} base model ---
chosen per-prompt by a lightweight complexity sensor and a Mamdani
fuzzy-logic controller --- rather than hosting multiple distinct models or
committing to one static bit-width. Unlike learned inter-model routers
(e.g., RouteLLM), this design requires only one set of base weights plus
small per-tier low-rank adapters resident in memory. We present the full
system design (complexity sensor, fuzzy inference system, precision tiers,
quantization-aware-trained LoRA adapters) with explicit mathematical
formalization, and a hardware-grounded measurement methodology built on
GPU on-board NVML energy counters together with an automated validation
gate that must pass before any number is eligible to be cited as a result.
This paper's empirical contribution at the present stage is a direct,
hardware-measured answer to the precondition the rest of the system
depends on: does energy-per-token actually differ across precision tiers
on real hardware, and by how much? Measuring Llama-3.2-1B on an NVIDIA RTX
6000 Ada Generation GPU across 4-bit, 8-bit, and 16-bit tiers
($n=1{,}500$ prompts per tier, 3 repeated runs, seed $=42$), we find mean
energy of $1.414\pm0.055$, $3.052\pm0.039$, and $8.124\pm0.241$ J/token
respectively, with non-overlapping 95\% confidence intervals across all
three tiers --- a real, monotonic, statistically separated energy
gradient across precision. This result passes this project's own
pre-registered go/no-go gate for whether dynamic precision routing is
worth building at all. The routing energy/accuracy tradeoff, baseline
comparisons against naive routers, and a quantization-aware-training
adapter ablation remain open empirical questions at the time of this
draft: the measurement sessions that will answer them are implemented and
scripted but not yet executed. We report their exact protocol so the
remaining results can be added without redesigning the system, and we are
explicit throughout about what is and is not yet known.
\end{abstract}

\begin{IEEEkeywords}
energy-proportional inference, dynamic quantization, fuzzy logic control,
LLM inference routing, green AI, hardware energy measurement
\end{IEEEkeywords}

\section{Introduction}
\label{sec:intro}

Serving LLM inference at a single, statically chosen numeric precision
treats every request identically regardless of how difficult it actually
is: a two-word factual query and a multi-step chain-of-reasoning math
problem pay the same per-token energy cost if both are served at, say,
16-bit. Prior work has separately explored (a) routing requests across
\emph{different models} by predicted preference (RouteLLM), (b) cascading
across \emph{paid API tiers} with a stop/escalate judger (FrugalGPT), and
(c) \emph{static} post-training or quantization-aware compression of a
single model (GPTQ, AWQ, LLM.int8, QLoRA). None of these directly asks
whether the numeric precision of a single, already-resident model can
itself be treated as a per-request routing decision --- chosen
automatically from an estimate of how complex a given prompt is, before
generation begins. That is the central design question of this paper's
system, Green-Weight: a lightweight, five-feature complexity sensor and
Mamdani fuzzy inference system route each prompt to one of three precision
tiers (4-bit, 8-bit, 16-bit) of a single base model, so only one set of
base weights plus small per-tier LoRA adapters must be resident, rather
than multiple hosted models.

Before any such system's energy claim can be trusted, however, its
underlying physical premise has to be checked on real hardware rather than
assumed: does energy-per-token actually fall meaningfully as bit-width
drops for this model class, or does GPU dequantize-to-fp16 compute (as it
does for some quantization kernels) erase the expected savings? This
project's own history is a caution here --- an earlier internal estimate
of end-to-end savings was derived from an assumed linear bit-width energy
model rather than measurement, and is explicitly retracted and not relied
upon anywhere in this paper (see \texttt{CREDIBILITY\_REPORT.md} \S 2 for
the full account and the project's standing rule against restating it).
This paper's methodological stance is the direct response to that
history: every claim below is either (i) traceable to a specific row of a
verifier-gated, append-only experiment log, or (ii) explicitly marked as
not yet measured. We do not report an aggregate ``energy savings'' or
``accuracy loss'' percentage anywhere in this paper, now or as a
projection, because no such number is licensed by measured data at the
time of writing.

Concretely, this paper reports:
\begin{enumerate}
  \item A formal system design for complexity-aware dynamic precision
    routing --- a five-feature complexity sensor (\S\ref{sec:sensor}) and
    a Mamdani fuzzy inference system with an explicit rule base,
    aggregation operator, and centroid defuzzification
    (\S\ref{sec:fuzzy}), both specified with numbered equations directly
    against the deployed implementation rather than described only in
    prose.
  \item A hardware-grounded measurement methodology (\S\ref{sec:method}):
    GPU on-board NVML energy counters (not a software or
    emissions-derived estimate), a fixed greedy-decoding protocol, and an
    automated validation gate (\texttt{verify\_results.py}) that every
    number in this paper must pass before it may be cited.
  \item The first hardware-measured result in this research program (RQ1,
    \S\ref{sec:rq1}): joules-per-token differs substantially and
    monotonically across the 4-bit / 8-bit / 16-bit precision tiers of
    Llama-3.2-1B on an NVIDIA RTX 6000 Ada Generation GPU, with
    non-overlapping 95\% confidence intervals across all three tiers.
    This is this project's pre-registered go/no-go gate for whether the
    rest of the routing system is worth evaluating, and it passes.
  \item A fully specified, not-yet-executed remainder of the empirical
    program --- routing energy/accuracy tradeoff (RQ2), baseline
    comparison against naive routers and an oracle (RQ3), and a
    QAT-adapter ablation (RQ4) --- with the exact GPU sessions that will
    produce each (\S\ref{sec:rq2}--\S\ref{sec:rq4}), so this paper's
    present evidentiary boundary is unambiguous to the reader.
\end{enumerate}

\textbf{Contributions:}
\begin{itemize}
  \item A complete mathematical formalization of the complexity sensor
    and Mamdani fuzzy routing controller (\S\ref{sec:system}), derived
    directly from the deployed implementation rather than approximated.
  \item A formal energy-accounting model (\S\ref{sec:energymodel}) that
    separates router compute overhead from tier inference cost, used as
    the basis for a break-even framing carried into future work
    (\S\ref{sec:conclusion}).
  \item An automated, pre-registered validation gate that must pass
    before any number is eligible to appear as a result in this paper
    (\S\ref{sec:gate}), and a full accounting of what has and has not yet
    passed it (\S\ref{sec:results}).
  \item A hardware-measured, statistically significant, monotonic energy
    separation across three precision tiers of a real 1B-parameter LLM on
    a commodity workstation GPU (\S\ref{sec:rq1}) --- the empirical
    premise the rest of this research program depends on.
  \item Honest, explicit scoping of the three research questions
    (RQ2--RQ4) that remain open at the time of this draft, with no
    projected or implied outcome for any of them
    (\S\ref{sec:rq2}--\S\ref{sec:rq4}, \S\ref{sec:conclusion}).
\end{itemize}

\section{Related Work}
\label{sec:related}

\textit{Bibliographic details (author lists, years, venues) below are
given to the best of the drafting agent's knowledge and should be checked
against primary sources before submission; none of this section depends
on this project's own measured numbers.}

\subsection{Learned Inter-Model Routing}

RouteLLM~\cite{routellm} trains a preference-predicting router that,
given a prompt, decides whether to send it to a cheap ``weak'' model or an
expensive ``strong'' model, learned from human-preference battle data
between specific named model pairs (e.g., GPT-4 vs.\ Mixtral).
Green-Weight addresses a related but structurally different problem:
rather than choosing between distinct models that must each be hosted, it
routes a single prompt across precision tiers of one base model
(\texttt{meta-llama/Llama-3.2-1B}, \S\ref{sec:tiers}), so only one set of
base weights plus small per-tier adapters (\S\ref{sec:adapters}) needs to
be resident. This sidesteps the multi-model hosting cost central to
RouteLLM's trade-off, but also means RouteLLM's own pretrained routers
cannot be reused as-is: its checkpoints are trained on preference battles
between specific named models drawn from a closed model-identity set and
carry no learned signal for ``should this prompt run at 4-bit, 8-bit, or
16-bit of the same model.'' Accordingly, the RouteLLM-shaped component of
this project's codebase (\texttt{router/routellm\_bridge.py}) is currently
a documented, provisional threshold pass-through over the fuzzy
controller's own output (\S\ref{sec:fuzzy}) --- it does not load or call a
RouteLLM model, and this paper makes no claim of a working RouteLLM
integration. A real tier-preference router, trained on this project's own
per-prompt routing data once that data exists, is planned future work and
is the actual point of comparison to RouteLLM's training methodology; it
is not part of the present contribution.

\subsection{Cascade-Based Cost Reduction}

FrugalGPT~\cite{frugalgpt} reduces the cost of querying paid third-party
LLM APIs by cascading: starting with a cheap model and using a learned
generation judger to decide whether an answer is good enough or whether
to escalate to a more expensive model. This project's codebase includes
an analogous path (\texttt{cascade/frugal\_cascade.py}) that wraps
FrugalGPT's own \texttt{LLMCascade} class around this project's local
precision tiers, but its escalation logic is not built out: the current
implementation performs a single call at the router-assigned starting
tier and returns a hardcoded judger score rather than actually invoking
FrugalGPT's judger and conditionally escalating to a higher tier. This
cascade path is treated as optional, deprioritized scope --- the paper's
core claim is complexity-aware precision routing via the fuzzy controller
(\S\ref{sec:fuzzy}), not the cascade --- and, if included at all, serves
at most as a secondary baseline, not the core contribution.

\subsection{Static and Post-Training Quantization}

GPTQ~\cite{gptq}, AWQ~\cite{awq}, LLM.int8()~\cite{llmint8}, and
QLoRA~\cite{qlora} address a different question: given a fixed precision
budget, how to quantize (or fine-tune under quantization) with minimal
quality loss. Each picks one static precision for an entire deployment.
This project's approach is complementary rather than competing: it uses
bitsandbytes' NF4 4-bit and int8 8-bit quantization --- the same family
of technique as LLM.int8()/QLoRA --- as two of its three precision tiers
(\S\ref{sec:tiers}), but treats precision itself as a per-request routing
decision rather than a fixed deployment-time choice, and trains a
separate QAT LoRA adapter for each tier (\S\ref{sec:adapters}) rather
than a single adapter tuned for one quantization level, applying QLoRA's
adapter-based recovery strategy independently at each of the three tiers.

\subsection{Adaptive Computation}

This work sits inside the broader adaptive-computation paradigm ---
early-exit networks, mixture-of-depths, speculative decoding, and other
methods that vary the amount of compute spent on a given input based on
its estimated difficulty. Where that literature typically adapts network
depth, number of decoding steps, or which draft/verify model pair is
used, this project adapts numeric precision: the same architecture and
layer depth run for every request, but the arithmetic precision of the
weights (and, correspondingly, memory-bandwidth demand) is selected per
request from a prompt-complexity estimate computed before generation
begins (\S\ref{sec:sensor}).

\subsection{Energy and Carbon Measurement Tooling}

Zeus~\cite{zeus} and LLMCarbon~\cite{llmcarbon} are representative of a
growing body of work on measuring and modeling the energy/carbon
footprint of deep learning workloads; MELODI-style GPU energy-profiling
tooling occupies similar territory (citation to be confirmed before
submission). This project follows their emphasis on grounding energy
claims in real measurement rather than an analytic model, but differs in
instrument: rather than software-level power estimation or
CO$_2$-emissions-derived accounting (an earlier version of this
project's own pipeline mis-derived energy by inverting CodeCarbon's
CO$_2$ output with an incorrect constant --- see
\texttt{CREDIBILITY\_REPORT.md} \S 2 --- a defect since fixed), this
project reads the GPU's onboard cumulative hardware energy counter
directly via NVML's \texttt{nvmlDeviceGetTotalEnergyConsumption}
(\S\ref{sec:method}), a hardware measurement rather than a software
estimate. This is a narrower but more directly verifiable instrument than
a full carbon-accounting framework: it does not model upstream grid
carbon intensity or embodied hardware carbon the way LLMCarbon does, and
it reports GPU-board joules per token only, explicitly excluding CPU/DRAM
host energy (\S\ref{sec:method}, \S\ref{sec:threats}).

\section{System}
\label{sec:system}

\subsection{Prompt-Complexity Sensor}
\label{sec:sensor}

The complexity sensor (\texttt{router/complexity\_scorer.py}) converts a
raw prompt string $p$ into five features that the fuzzy controller
(\S\ref{sec:fuzzy}) consumes, four of them normalized to $[0,1]$ and one
binary:

\begin{enumerate}
  \item \textbf{Flesch--Kincaid grade level} (via
    \texttt{textstat.flesch\_kincaid\_grade}), normalized against
    \texttt{FLESCH\_KINCAID\_RANGE = (0, 18)}.
  \item \textbf{Token length}, approximated as prompt character count
    divided by 4 (a standard English characters-per-token heuristic),
    normalized against \texttt{TOKEN\_LENGTH\_RANGE = (0, 154)}.
  \item \textbf{Character-level Shannon entropy}, normalized against
    \texttt{ENTROPY\_RANGE = (3.3, 5.0)} bits.
  \item \textbf{Syntax-tree depth}, the maximum dependency-parse depth
    returned by spaCy's \texttt{en\_core\_web\_sm} model, normalized
    against \texttt{SYNTAX\_DEPTH\_RANGE = (2, 14)}.
  \item \texttt{has\_code\_or\_math}, a binary (unnormalized) feature set
    by regex matching against markdown/fenced code blocks, inline code,
    LaTeX delimiters and commands, common mathematical operators
    ($\sum$, $\int$, $\partial$, $\sqrt{\cdot}$, $\infty$, and relational
    symbols), and a keyword list of natural-language code/algorithm
    requests (e.g.\ ``implement,'' ``algorithm,'' ``recursion'').
\end{enumerate}

All continuous features share a single \texttt{normalize\_to\_01(value,
min, max)} helper that clips to $[0,1]$; the fuzzy controller
(\S\ref{sec:fuzzy}) reuses the same range constants when normalizing its
own membership-function breakpoints, so the two stay aligned by
construction rather than by convention.

\textbf{Formal definition.} Let the normalization operator be
\begin{equation}
\text{norm}(x; a, b) = \mathrm{clip}\!\left(\frac{x - a}{b - a},\ 0,\ 1\right)
\label{eq:norm}
\end{equation}
used identically for every continuous feature below. The sensor produces
a 5-dimensional feature vector for prompt $p$:
\begin{equation}
f(p) = [f_{FK},\ f_{TL},\ f_{H},\ f_{SD},\ f_{CM}]
\label{eq:featurevec}
\end{equation}
where
\begin{equation}
f_{FK} = \text{norm}(FK(p);\ 0,\ 18)
\label{eq:ffk}
\end{equation}
with $FK(p)$ the Flesch--Kincaid grade level of $p$;
\begin{equation}
f_{TL} = \text{norm}\!\left(\left\lfloor \frac{\mathrm{len}(p)}{4} \right\rfloor;\ 0,\ 154\right)
\label{eq:ftl}
\end{equation}
\begin{equation}
f_{H} = \text{norm}(H(p);\ 3.3,\ 5.0),\quad H(p) = -\sum_{c \in \Sigma} p(c)\log_2 p(c)
\label{eq:fh}
\end{equation}
with $H(p)$ the character-level Shannon entropy over the empirical
character distribution of $p$ ($\Sigma$ = the set of distinct characters
appearing in $p$);
\begin{equation}
f_{SD} = \text{norm}(D(p);\ 2,\ 14)
\label{eq:fsd}
\end{equation}
with $D(p)$ the maximum depth of the spaCy dependency parse of $p$; and
\begin{equation}
f_{CM} \in \{0, 1\},\ \text{an indicator over the regex/keyword match above.}
\label{eq:fcm}
\end{equation}
$f_{CM}$ is explicitly \emph{not} passed through $\text{norm}(\cdot)$ ---
it is a binary indicator, not a continuous normalized feature.

The range constants in Eqs.~(\ref{eq:ffk}), (\ref{eq:fh}),
(\ref{eq:fsd}), (\ref{eq:ftl}) --- $(0,18)$, $(3.3,5.0)$, $(2,14)$,
$(0,154)$ --- are not arbitrary defaults; they are empirically calibrated
against this project's real 500-prompt stratified evaluation set (200
easy / 150 medium / 150 hard prompts, seed $=42$,
\texttt{data/eval\_prompts.jsonl}), following an earlier 30-prompt
calibration pass. The original entropy and syntax-depth ranges clipped
the real distribution's upper tail --- syntax depth in particular
exceeded the old bound for roughly the hardest 5\% of prompts, collapsing
them to an identical flat 1.0 regardless of their actual relative
difficulty --- and the original token-length normalization used a
hardcoded 512-token cap that the real prompt set never approached,
leaving that feature structurally unable to signal difficulty in
practice. \texttt{TOKEN\_LENGTH\_RANGE = (0, 154)} was derived directly
from the observed distribution of the real evaluation set. The values in
Eqs.~(\ref{eq:ffk})--(\ref{eq:fsd}) are the corrected, re-calibrated
constants currently in use (see \texttt{CREDIBILITY\_REPORT.md}, row
``Fuzzy controller's feature-normalization bounds,'' for the full
calibration history).

One caveat carried into the routing design rather than resolved by range
calibration alone: an early per-difficulty analysis of this dataset
suggested character-level entropy ($f_H$) discriminates prompt difficulty
only weakly, while syntax depth ($f_{SD}$) showed clearer separation
between easy and hard prompts. This motivates treating entropy as a
candidate for a feature-ablation study rather than assuming all five
features contribute equally to routing quality; we flag it here as a
design consideration, and defer any quantitative ablation result to
\S\ref{sec:rq4} (not yet measured).

\subsection{Fuzzy Gearbox Controller}
\label{sec:fuzzy}

The routing decision is made by a Mamdani-type fuzzy inference system
(\texttt{router/fuzzy\_controller.py}), implemented with the
\texttt{scikit-fuzzy} library rather than a hand-rolled approximation of
one. Each of the five complexity-sensor features becomes a fuzzy
antecedent over the universe $[0,1]$: \texttt{flesch\_kincaid},
\texttt{token\_length}, \texttt{entropy}, and \texttt{syntax\_depth} each
get three triangular membership functions --- LOW, MEDIUM, HIGH --- built
from two edge breakpoints read from \texttt{config.yaml} (LOW spans
$[0,0,lo]$, MEDIUM spans $[lo, mid, hi]$ with $mid$ as their midpoint,
HIGH spans $[hi,1,1]$). Breakpoints are stored in each feature's raw
units (e.g.\ entropy bits, parse-tree depth) and normalized through the
\emph{same} range constants as \S\ref{sec:sensor}
(Eqs.~(\ref{eq:ffk}), (\ref{eq:fh}), (\ref{eq:fsd})) before being used as
fuzzy-set edges, so the membership functions stay aligned with the
features they are compared against --- a real design property enforced
in code, not an incidental convenience. \texttt{has\_code\_or\_math} gets
two membership functions (``no''/``yes'') split at 0.5, matching its
effectively binary nature.

\textbf{Formal definition.} The single membership-function primitive used
throughout is the triangular function
\begin{equation}
\mathrm{tri}(x; a, b, c) = \max\!\left(\min\!\left(\frac{x-a}{b-a},\ \frac{c-x}{c-b}\right),\ 0\right)
\label{eq:tri}
\end{equation}
For each continuous antecedent feature $x \in \{f_{FK}, f_{TL}, f_H,
f_{SD}\}$ with calibrated breakpoints $(lo, hi)$ read from
\texttt{config.yaml} and normalized into the same $[0,1]$ space as $x$
via Eq.~(\ref{eq:norm}), and $mid = (lo+hi)/2$:
\begin{align}
\mu_{LOW}(x) &= \mathrm{tri}(x;\ 0,\ 0,\ lo) \label{eq:mulow}\\
\mu_{MEDIUM}(x) &= \mathrm{tri}(x;\ lo,\ mid,\ hi) \label{eq:mumed}\\
\mu_{HIGH}(x) &= \mathrm{tri}(x;\ hi,\ 1,\ 1) \label{eq:muhigh}
\end{align}
For the binary antecedent $f_{CM}$:
\begin{align}
\mu_{NO}(f_{CM}) &= \mathrm{tri}(f_{CM};\ 0,\ 0,\ 0.5) \label{eq:muno}\\
\mu_{YES}(f_{CM}) &= \mathrm{tri}(f_{CM};\ 0.5,\ 1,\ 1) \label{eq:muyes}
\end{align}
The consequent ``complexity'' $y \in [0,100]$ has three fixed membership
functions, matching the deployed configuration exactly:
\begin{align}
\mu_{Z=LOW}(y) &= \mathrm{tri}(y;\ 0,\ 0,\ 33) \label{eq:zlow}\\
\mu_{Z=MEDIUM}(y) &= \mathrm{tri}(y;\ 25,\ 50,\ 75) \label{eq:zmed}\\
\mu_{Z=HIGH}(y) &= \mathrm{tri}(y;\ 67,\ 100,\ 100) \label{eq:zhigh}
\end{align}

The rule base consists of seven Mamdani IF-THEN rules $R_1,\dots,R_7$:
$R_1$ --- low readability, low token length, low syntax depth, and no
code/math $\Rightarrow$ LOW; $R_2$ --- (high syntax depth OR high
entropy) AND high readability grade $\Rightarrow$ HIGH; $R_3$ --- code or
math detected $\Rightarrow$ HIGH; $R_4$ --- high token length
$\Rightarrow$ MEDIUM; $R_5$ --- high entropy AND high readability grade
$\Rightarrow$ HIGH; $R_6$ --- medium token length OR medium syntax depth
$\Rightarrow$ MEDIUM; $R_7$ --- low readability AND low token length
$\Rightarrow$ LOW. For each rule $R_k$, its firing strength combines
antecedent memberships via the standard Mamdani/Zadeh operators --- $\min$
for AND ($\wedge$), $\max$ for OR ($\vee$):
\begin{equation}
\alpha_k = T\big(\mu_{A_{i_1}}(x_{i_1}), \dots, \mu_{A_{i_n}}(x_{i_n})\big)
\label{eq:firing}
\end{equation}
Mamdani implication clips each rule's consequent membership function at
its firing strength, and rule outputs are aggregated across all seven
rules via $\max$:
\begin{equation}
\mu_{agg}(y) = \max_k \min\big(\alpha_k,\ \mu_{Z_k}(y)\big)
\label{eq:agg}
\end{equation}
Centroid (center-of-gravity) defuzzification, computed by
\texttt{scikit-fuzzy}'s own library implementation rather than custom
project code, produces the crisp complexity score:
\begin{equation}
y^{*} = \frac{\int y \cdot \mu_{agg}(y)\, dy}{\int \mu_{agg}(y)\, dy}
\label{eq:centroid}
\end{equation}
The crisp score maps to a precision tier via configurable cut points
(default 33/66):
\begin{equation}
\text{tier}(y^{*}) =
\begin{cases}
\text{4-bit} & y^{*} \le 33 \\
\text{8-bit} & 33 < y^{*} \le 66 \\
\text{16-bit} & y^{*} > 66
\end{cases}
\label{eq:tiermap}
\end{equation}
and
\begin{equation}
\text{win\_probability} = y^{*}/100
\label{eq:winprob}
\end{equation}
Despite the name --- kept for interface compatibility with the
RouteLLM-shaped bridge described in \S\ref{sec:related}A --- Eq.~(\ref{eq:winprob})
is \emph{not} a RouteLLM classifier's output; it is a deterministic
rescaling of this system's own fuzzy defuzzification result, consumed by
\texttt{router/routellm\_bridge.py}'s threshold pass-through to make the
final tier decision.

\subsection{Precision Tiers}
\label{sec:tiers}

Each prompt is routed to one of three precision tiers of a single base
model, \texttt{meta-llama/Llama-3.2-1B}:
\begin{itemize}
  \item \textbf{4-bit}: bitsandbytes NF4 quantization with double
    quantization enabled (\texttt{BitsAndBytesConfig(load\_in\_4bit=True,
    bnb\_4bit\_quant\_type="nf4", bnb\_4bit\_use\_double\_quant=True)}).
  \item \textbf{8-bit}: bitsandbytes LLM.int8() quantization
    (\texttt{load\_in\_8bit=True}).
  \item \textbf{16-bit}: bfloat16, loaded without quantization as the
    full-precision reference tier.
\end{itemize}
During measurement (\S\ref{sec:method}), only one tier is held in GPU
memory at a time, so that each tier's energy and latency behavior is
measured in isolation rather than under any shared-memory or
co-residency effect.

\subsection{QAT LoRA Adapters}
\label{sec:adapters}

Each precision tier is paired with its own quantization-aware-training
(QAT) LoRA adapter, trained independently rather than sharing a single
adapter across tiers: \texttt{adapter\_simple} for the 4-bit tier (rank
$r=16$, $\alpha=32$), \texttt{adapter\_medium} for the 8-bit tier
($r=8$, $\alpha=16$), and \texttt{adapter\_complex} for the 16-bit tier
($r=4$, $\alpha=8$) --- a constant $\alpha/r$ ratio of 2 across all
three, with LoRA rank scaling inversely with precision. The design
rationale is that the most aggressively quantized tier (4-bit) has the
most quantization error to recover from and is given the most adapter
capacity to do so, while the least-quantized tier (16-bit) needs the
least. All three adapters target the same two attention projections
(\texttt{q\_proj}, \texttt{v\_proj}) and share the
\texttt{meta-llama/Llama-3.2-1B} base.

The three adapters' existence, base-model consistency, and loadability
have been confirmed by a real, non-simulated load test: the actual gated
\texttt{meta-llama/Llama-3.2-1B} weights were downloaded, each adapter
was attached via \texttt{PeftModel.from\_pretrained()}, and each was
exercised with a real greedy-decoding generation, with no errors across
all three. This confirms the adapters are usable, correctly-shaped
artifacts; it does not by itself establish that they improve accuracy
over plain post-training quantization at each tier --- that comparison is
a measurement question (per-tier accuracy evaluation, not yet complete as
of this draft; see \S\ref{sec:rq4}) and is not claimed here.

\section{Measurement Methodology}
\label{sec:method}

\subsection{Energy Measurement}

Energy is measured with the GPU's onboard hardware energy counter, read
via NVML's \texttt{nvmlDeviceGetTotalEnergyConsumption} (millijoule
resolution) --- a direct hardware measurement, not a software estimate or
an emissions-derived proxy (see \S\ref{sec:related}E's note on this
project's earlier, now-fixed CO$_2$-inversion bug). If a node's driver
does not expose the counter, the measurement script falls back to
integrating 20\,Hz power samples and the affected run is flagged for
disclosure rather than reported as equivalent to a counter-based
measurement.

The protocol for each tier is: 5 warmup generations (discarded from
measurement), followed by greedy (deterministic) decoding over the
evaluation prompt set with a fixed \texttt{max\_new\_tokens=128} and real
generated token counts (not an estimate). Only one precision tier is
resident in GPU memory at a time (\S\ref{sec:tiers}).
\texttt{torch.cuda.synchronize()} brackets each measured generation to
avoid attributing asynchronous CUDA work to the wrong measurement window.
An idle-power baseline and the GPU/driver identity (model, driver
version, clocks) are recorded to \texttt{hardware\_info.json} for every
run. The full sweep across tiers is repeated three times on different
days, and results are reported as mean $\pm$ 95\% confidence interval.

Because decoding is greedy and therefore deterministic, a router's
per-prompt output and energy cost are identical to those of whichever
tier it selects for that prompt; routing-condition results are
accordingly derived from the same per-tier, per-prompt measurement grid
rather than re-run independently.

\subsection{Energy Accounting Model}
\label{sec:energymodel}

We formalize the per-request energy accounting used throughout this
paper and its remaining measurement program. For a prompt $p$ routed to
tier $\text{tier}(p)$ (Eq.~(\ref{eq:tiermap})), total energy is
\begin{equation}
E_{total}(p) = E_{router}(p) + E_{inference}(\text{tier}(p),\ p)
\label{eq:etotal}
\end{equation}
where $E_{router}(p)$ is the fuzzy controller's own CPU-side compute cost
(feature extraction plus fuzzy inference) and $E_{inference}$ is the GPU
energy consumed generating the response at the chosen tier. Per-tier
mean joules-per-token, exactly what \S\ref{sec:rq1} reports, is
\begin{equation}
\bar{e}_{tier} = \frac{1}{N}\sum_{i=1}^{N} \frac{E_i}{T_i},\quad
\text{reported as } \bar{e}_{tier} \pm 1.96\frac{s}{\sqrt{N}}\ \text{(95\% CI)}
\label{eq:ebar}
\end{equation}
with $N = 1{,}500$ per tier in the measurement reported in
\S\ref{sec:rq1}. Eq.~(\ref{eq:etotal}) makes explicit that $E_{router}$
is measured and reported separately rather than assumed to be zero,
since it runs on CPU in milliseconds and is not free: quantifying it is
necessary to determine whether, and at what prompt lengths, routing
overhead is worth paying relative to always serving at a single static
tier (the break-even framing carried into \S\ref{sec:conclusion} as
future work). As of this draft, $E_{router}$ has not yet been measured on
the college GPU cluster and no number is reported for it; this is a
planned, required measurement (project roadmap, Phase 6), not an
assumption of zero cost.

\subsection{Accuracy Measurement}

Per-tier benchmark accuracy uses \texttt{lm-eval-harness}, evaluating
community-standard, versioned tasks with a fixed seed (42): tinyMMLU, a
GSM8K subset, and a HellaSwag subset, chosen to fit the measurement
budget while remaining standard, citable benchmarks. Raw per-task JSON
output is retained.

A separate, explicitly-labeled proxy is used to score correctness on
this project's own routed evaluation set (as distinct from the standard
lm-eval-harness benchmark tasks above): a reference-match check ---
numeric match for arithmetic/math answers, normalized substring match
for short factual answers, and token-level F1 $\ge 0.5$ otherwise. This
proxy is used only to attribute per-prompt correctness to a routing
decision on this project's own prompt set; it is not a substitute for,
and is kept explicitly distinct from, the lm-eval-harness benchmark
accuracy numbers, which remain this paper's accuracy claim of record.

\subsection{Automated Validation Gate}
\label{sec:gate}

Every results file produced by a measurement session is checked by
\texttt{training/scripts/verify\_results.py} before any number from it
may be cited anywhere in this paper. The gate checks: the output is
non-empty; each tier has at least $n=30$ samples; all recorded energies
and token counts are positive; joules-per-token values fall inside a
physical plausibility window for a roughly 1-billion-parameter model
(0.05--50 J/token); confidence intervals are acceptably tight; at least 3
repeated runs exist; and a set of internal-consistency checks hold
(e.g., an oracle router's accuracy must be at least as high as any
non-oracle condition's; a result implying 4-bit inference costs more
energy than 16-bit would trigger a disclosed risk warning rather than
being silently accepted). Only rows in \texttt{paper/results.md} with a
passing verdict from this gate are eligible to be cited as results in
this paper; any WARN is either resolved before citing the affected
number or explicitly carried into \S\ref{sec:threats} (Threats to
Validity).

\subsection{Current Status}

All measurements described above run on a college GPU cluster (an
NVIDIA RTX 6000 Ada Generation card), not the T4 originally planned for;
this platform switch is documented in the project roadmap. As of this
draft, per \texttt{paper/results.md}, the energy-measurement protocol
above (Session 1) has been executed as a full sweep over the 500-prompt
evaluation set, repeated three times, and has passed the automated
validation gate described above; the resulting per-tier energy figures
are reported in \S\ref{sec:rq1}, not here, since this section is
confined to methodology. Per-tier accuracy evaluation (Session 2, via
lm-eval-harness, \S\ref{sec:method}C) and the main routing experiment
(Session 4) had not, as of this draft, produced a validator-passing,
citable result --- no accuracy or routing numbers are stated anywhere in
this paper until they do.

\section{Results}
\label{sec:results}

\subsection{RQ1 (Measurement) --- Answered}
\label{sec:rq1}

\textbf{Finding.} Energy per token differs substantially and
monotonically across the three precision tiers of Llama-3.2-1B on an
NVIDIA RTX 6000 Ada Generation GPU. Table~\ref{tab:energy} reports the
full result, drawn verbatim from Session 1 of \texttt{paper/results.md}
(2026-09-04, SPIT cluster; driver 610.43.02, idle power 22.53\,W;
500-prompt evaluation set, 3 full sweeps, seed $=42$).

\begin{table}[h]
\caption{Measured energy per token by precision tier (Eq.~(\ref{eq:ebar})).}
\label{tab:energy}
\centering
\begin{tabular}{@{}lrrl@{}}
\toprule
Tier & $n$ & Mean J/token & 95\% CI bounds \\
\midrule
4-bit  & 1,500 & $1.414 \pm 0.055$ & $[1.359,\ 1.469]$ \\
8-bit  & 1,500 & $3.052 \pm 0.039$ & $[3.013,\ 3.091]$ \\
16-bit & 1,500 & $8.124 \pm 0.241$ & $[7.883,\ 8.365]$ \\
\bottomrule
\end{tabular}
\end{table}

Pairwise ratios computed from these means:
\begin{equation}
\frac{\bar{e}_{8bit}}{\bar{e}_{4bit}} = \frac{3.052}{1.414} \approx 2.16\times,\quad
\frac{\bar{e}_{16bit}}{\bar{e}_{8bit}} = \frac{8.124}{3.052} \approx 2.66\times,\quad
\frac{\bar{e}_{16bit}}{\bar{e}_{4bit}} = \frac{8.124}{1.414} \approx 5.75\times
\label{eq:ratios}
\end{equation}

The three tiers' 95\% confidence intervals do not overlap at any adjacent
pair: 4-bit's upper bound (1.469) is below 8-bit's lower bound (3.013),
and 8-bit's upper bound (3.091) is below 16-bit's lower bound (7.883).
This is the statistical basis for treating the separation as real rather
than measurement noise: energy-per-token increases monotonically and
significantly from 4-bit to 8-bit to 16-bit.

\textbf{Interpretation against the go/no-go gate.} The project roadmap
defines this project's go/no-go criterion as: ``if 4-bit/8-bit show real,
non-trivial energy savings over fp16 $\rightarrow$ proceed.'' The
measured, non-overlapping, monotonic separation above satisfies that
criterion on real hardware --- this is not the assumed linear bit-width
model referenced in \S\ref{sec:intro} as retracted, but a directly
measured hardware result. This licenses proceeding with the routing
system's remaining evaluation (RQ2--RQ4 below): the physical premise that
precision tier materially changes per-token energy cost on this hardware
is empirically established.

\textbf{Scope.} This result is a single-GPU-class (NVIDIA RTX 6000 Ada
Generation), single-model (Llama-3.2-1B) measurement. We do not claim it
generalizes to other GPU architectures, other model sizes/families, or
other decoding configurations; see \S\ref{sec:threats} for the full
scope discussion.

\subsection{RQ2 (Routing) --- Pending}
\label{sec:rq2}

RQ2 asks how much energy complexity-aware fuzzy routing
(\S\ref{sec:fuzzy}) uses relative to static fp16 inference, and at what
accuracy cost, across the full 500-prompt evaluation set. This requires
running all routing conditions (static tiers, fuzzy router, and the
baseline/oracle conditions listed under RQ3 below) over the same
prompt$\times$tier measurement grid described in \S\ref{sec:method},
produced by Session 4 (\texttt{training/scripts/kaggle\_routing\_}
\texttt{experiment.py}). This session is implemented and scripted but has
not yet been run on the GPU cluster as of this draft. No energy,
accuracy, or energy/accuracy-tradeoff number for routing is stated
anywhere in this paper; this subsection will be filled in from a
validator-passing row of \texttt{paper/results.md} once Session 4
completes.

\subsection{RQ3 (Baselines) --- Pending}
\label{sec:rq3}

RQ3 asks whether the fuzzy controller (\S\ref{sec:fuzzy}) outperforms
trivial routing strategies --- a random router with tier distribution
matched to the fuzzy router, and a naive threshold router over the mean
of the five raw complexity features (\texttt{naive\_complexity\_score()},
deliberately built as a non-tautological baseline distinct from the
fuzzy controller's own defuzzified output) --- and how close it
approaches an oracle router (cheapest tier that still answers correctly).
This also requires Session 4's routing experiment, evaluated over the
same conditions as RQ2. No baseline-comparison result is stated anywhere
in this paper as of this draft.

\subsection{RQ4 (Ablation) --- Pending}
\label{sec:rq4}

RQ4 asks whether the per-tier QAT LoRA adapters (\S\ref{sec:adapters})
improve accuracy at each precision tier relative to plain post-training
quantization without an adapter. This requires Session 2
(\texttt{training/scripts/kaggle\_accuracy\_eval.py}), which evaluates
each tier with and without its adapter via lm-eval-harness
(\S\ref{sec:method}C). Session 2 is implemented and scripted but has not
yet been run on the GPU cluster as of this draft. No accuracy or
ablation number is stated anywhere in this paper; \S\ref{sec:adapters}'s
adapter load-test result (adapters attach and generate without error) is
a loadability check only and is not an accuracy claim.

\section{Threats to Validity}
\label{sec:threats}

Drawn from \texttt{CREDIBILITY\_REPORT.md} \S 5, restated against the
RQ1 result reported in \S\ref{sec:rq1} and this paper's known
measurement artifacts.

\begin{enumerate}
  \item \textbf{Hardware scope.} All energy measurements reported in
    this paper are from one GPU model, the NVIDIA RTX 6000 Ada
    Generation (\S\ref{sec:method}E, a platform switch from the T4
    originally planned). Findings may differ on other GPU architectures
    (e.g., A100, H100, or older Turing-class cards) or on CPU inference.
    Claims in this paper are scoped to this specific card and are not
    asserted to generalize across GPU classes without further
    measurement.
  \item \textbf{Model scope.} All results are from one model,
    Llama-3.2-1B. No claim is made about other model sizes or families
    without measurement.
  \item \textbf{Quantization-kernel reality vs.\ theoretical bit-width
    scaling.} bitsandbytes' 4-bit and 8-bit kernels dequantize weights
    to a higher precision for compute; energy is not guaranteed to
    scale linearly, or even monotonically, with nominal bit-width for
    every model/hardware combination. \S\ref{sec:rq1}'s result confirms
    a real, monotonic, statistically significant separation on this
    hardware for this model, but this is reported as a \emph{measured}
    outcome specific to this configuration, not as evidence that
    bit-width and energy scale proportionally in general.
  \item \textbf{Correctness proxy.} The reference-match proxy used to
    score correctness on this project's own routed evaluation set
    (\S\ref{sec:method}C) --- exact numeric match, normalized substring
    match, or token-F1 $\ge 0.5$ --- is an approximation of true answer
    correctness, not a substitute for lm-eval-harness's standard
    benchmark scoring. Any oracle-router condition (RQ3,
    \S\ref{sec:rq3}) that depends on this proxy inherits its noise.
  \item \textbf{Shared-infrastructure noise.} The college GPU cluster
    used for measurement is shared infrastructure; thermal and host
    conditions can vary between sessions. This is mitigated, not
    eliminated, by repeating the full sweep three times on different
    days and reporting 95\% confidence intervals (Eq.~(\ref{eq:ebar}));
    the non-overlapping CIs in Table~\ref{tab:energy} are evidence the
    observed tier separation exceeds this noise floor, but
    session-to-session variation beyond what three repeats capture
    cannot be ruled out.
  \item \textbf{Energy counter scope (GPU-board-only).} NVML's
    \texttt{nvmlDeviceGetTotalEnergyConsumption} measures GPU-board
    energy only; host CPU and DRAM energy are excluded from every
    number in this paper. Total system energy for a request is
    therefore higher than the reported J/token figures, which should be
    read as GPU-board energy specifically, not end-to-end system
    energy.
  \item \textbf{NVML counter resolution at very short generations.} 117
    of 4,500 measured rows (2.6\%, across all three tiers) recorded
    $\text{energy\_j} = 0.0$. These come from prompts that generated
    only a single output token, where the elapsed generation window is
    short enough that the NVML counter's millijoule-scale rounding
    resolution can report zero net energy change over that window. This
    is a real measurement artifact of counter resolution at short
    durations, not a meter failure or a discarded/invalid reading, and
    it is disclosed here as a limitation of the instrument at very
    short generations rather than silently excluded from the reported
    means; \texttt{verify\_results.py} flags but does not reject these
    rows, and Table~\ref{tab:energy}'s means are computed over the full
    1,500-row per-tier set inclusive of them.
\end{enumerate}

\section{Conclusion}
\label{sec:conclusion}

This paper reports the design and initial hardware-measured evaluation
of Green-Weight, a system that routes individual LLM requests across
precision tiers of a single base model using a five-feature complexity
sensor and a Mamdani fuzzy inference controller, formalized in
\S\ref{sec:system} with numbered equations against the deployed
implementation rather than described only in prose. We built a
measurement methodology grounded in GPU on-board hardware energy
counters and an automated, pre-registered validation gate
(\S\ref{sec:method}) so that no number in this paper is asserted without
a machine-checked, traceable source.

At the time of this draft, one of this paper's four research questions
is fully answered: RQ1 establishes, on real hardware, that
energy-per-token differs substantially and monotonically across the
4-bit, 8-bit, and 16-bit precision tiers of Llama-3.2-1B on an NVIDIA RTX
6000 Ada Generation GPU, with statistically non-overlapping confidence
intervals across all three tiers (\S\ref{sec:rq1}). This is this
project's own pre-registered go/no-go criterion for whether dynamic
precision routing is a physically grounded idea worth building the rest
of the system around, and it passes. RQ2 (routing energy/accuracy
tradeoff), RQ3 (comparison against naive and oracle baseline routers),
and RQ4 (QAT adapter ablation) remain open: the GPU sessions that will
answer them (Session 4 for RQ2/RQ3, Session 2 for RQ4) are implemented
and scripted against the same verifier-gated methodology used for RQ1,
but have not yet been executed as of this draft. We report their exact
protocol in \S\ref{sec:method} and \S\ref{sec:results} rather than any
projected outcome, so this paper's evidentiary boundary is unambiguous.

Near-term future work is therefore concrete rather than speculative:
running Sessions 2 and 4 to answer RQ2--RQ4; measuring the router's own
compute overhead, $E_{router}$ in Eq.~(\ref{eq:etotal}), which is
currently reported as an un-measured but explicitly non-zero quantity, to
complete the energy-accounting/break-even model sketched in
\S\ref{sec:energymodel} (the prompt-length point below which routing
overhead outweighs its tier-selection savings); training a real
tier-preference router on Session 4's per-prompt routing data as the
genuine RouteLLM-methodology comparison point discussed in
\S\ref{sec:related}A, superseding the current documented threshold
pass-through; and a memory-footprint comparison of this project's
single-base-model-plus-adapters design against the cost of hosting
multiple separate models under a RouteLLM-style routing scheme. We view a
paper that honestly reports one fully measured, statistically
significant finding alongside a fully specified but not-yet-executed
remainder as the appropriate way to report work at this project's
current stage, and we commit to reporting whatever RQ2--RQ4 show once
measured, whether or not the resulting numbers are as favorable as an
earlier, retracted, non-measured estimate once assumed.

% ============================================================================
% Bibliography — placeholder entries. Verify every citation against the
% primary source before submission; do not trust years/venues below blindly.
% ============================================================================
\begin{thebibliography}{00}

\bibitem{routellm} I.~Ong, A.~Almahairi, V.~Wu, W.-L.~Chiang, T.~Wu,
J.~E.~Gonzalez, M.~W.~Kadous, and I.~Stoica, ``RouteLLM: Learning to
Route LLMs with Preference Data,'' \emph{arXiv preprint}, 2024.

\bibitem{frugalgpt} L.~Chen, M.~Zaharia, and J.~Zou, ``FrugalGPT: How to
Use Large Language Models While Reducing Cost and Improving
Performance,'' \emph{arXiv preprint arXiv:2305.05176}, 2023.

\bibitem{gptq} E.~Frantar, S.~Ashkboos, T.~Hoefler, and D.~Alistarh,
``GPTQ: Accurate Post-Training Quantization for Generative
Pre-trained Transformers,'' \emph{arXiv preprint arXiv:2210.17323},
2022.

\bibitem{awq} J.~Lin, J.~Tang, H.~Tang, S.~Yang, X.~Dang, and S.~Han,
``AWQ: Activation-aware Weight Quantization for LLM Compression and
Acceleration,'' \emph{arXiv preprint arXiv:2306.00978}, 2023.

\bibitem{llmint8} T.~Dettmers, M.~Lewis, Y.~Belkada, and L.~Zettlemoyer,
``LLM.int8(): 8-bit Matrix Multiplication for Transformers at Scale,''
\emph{Advances in Neural Information Processing Systems}, 2022.

\bibitem{qlora} T.~Dettmers, A.~Pagnoni, A.~Holtzman, and L.~Zettlemoyer,
``QLoRA: Efficient Finetuning of Quantized LLMs,'' \emph{Advances in
Neural Information Processing Systems}, 2023.

\bibitem{zeus} J.~W.~Chung, J.~Kim, Z.~Zhang, M.~Song, N.~Kwon, B.~Chun,
and M.~Chowdhury, ``Zeus: Understanding and Optimizing GPU Energy
Consumption of DNN Training,'' \emph{USENIX NSDI}, 2023.

\bibitem{llmcarbon} A.~Faiz, S.~Kaneda, R.~Wang, R.~Osi, P.~Sharma,
F.~Chen, and L.~Jiang, ``LLMCarbon: Modeling the End-to-End Carbon
Footprint of Large Language Models,'' \emph{arXiv preprint
arXiv:2309.14393}, 2024.

\end{thebibliography}

\end{document}
